\documentclass[11pt]{article}

\usepackage{amsmath,amsfonts,amssymb,amsthm}
\usepackage{mathtools}
\usepackage{hyperref}
\usepackage{enumitem}
\usepackage{booktabs}
\usepackage{geometry}
\geometry{margin=1in}

% Common macros (keep small; avoid redefining paper content)
\newcommand{\R}{\mathbb{R}}
\newcommand{\N}{\mathbb{N}}
\newcommand{\E}{\mathbb{E}}
\newcommand{\Prb}{\mathbb{P}}
\newcommand{\supp}{\mathrm{supp}}
\newcommand{\TV}{\mathrm{TV}}
\newcommand{\KL}{\mathrm{D_{KL}}}
\newcommand{\Dinfty}{\mathrm{D_{\infty}}}
\newcommand{\MaxL}{\mathcal{L}}
\newcommand{\PML}{\mathcal{L}^{\bullet}}
\newcommand{\given}{\,|\,}
\newcommand{\Unif}{\mathrm{Unif}}
\newcommand{\eps}{\varepsilon}
\newcommand{\del}{\delta}
\newcommand{\argmax}{\operatorname*{arg\,max}}
\newcommand{\argmin}{\operatorname*{arg\,min}}
\newcommand{\bits}{\mathrm{bits}}
\newcommand{\logtwo}{\log_2}

% Sets / alphabets
\newcommand{\C}{\mathcal{C}} % secrets (destinations, profiles, committees)
\newcommand{\Y}{\mathcal{Y}} % observations
\newcommand{\W}{\mathcal{W}} % witnesses

% Convenience
\newcommand{\ind}{\mathbf{1}}


\title{Boss Fight Budgets:\\ Maximal Leakage as a Repeated-Use Anonymity Contract}
\author{}
\date{February 2026 (archive v125)}

\newtheorem{theorem}{Theorem}
\newtheorem{lemma}{Lemma}
\newtheorem{proposition}{Proposition}
\newtheorem{corollary}{Corollary}
\newtheorem{definition}{Definition}

\begin{document}
\maketitle

\begin{abstract}
The \emph{Boss Fight} is the repeated-use regime for metadata-hiding overlays: a determined analyst aggregates many protocol instances to identify a hidden secret (destination, key, profile, committee, or flow).
We argue that \emph{maximal leakage} (MaxL) and \emph{pointwise maximal leakage} (PML) are the right \emph{identification} budgets: they equal worst-case multiplicative gain in guessing and they compose under adaptivity.
We present (i) a conservative, audit-friendly step bound for common overlay patterns and (ii) a minimal calculation showing why ``target + dummies'' can \emph{collapse} under partial observation and selective failure, motivating higher-order witness constraints (``KICS'').
This paper is intentionally brief and delegates padding/scheduling mechanisms to earlier papers in this series.
\end{abstract}

\paragraph{Notation.}
We follow the minimal model and naming conventions in the series notation sheet (Synthesis~8): hidden index $K$, state $S$, transcript $T$, and declared projections $X=f(T)$.\cite{series:notation}

\paragraph{Scope.}
We formalize the \,\emph{Boss Fight}\, regime: repeated, adaptive use of an anonymity mechanism with an analyst optimizing identification.
The objective is an audit-friendly end-to-end \emph{identification} budget rather than a single-run indistinguishability claim.



\begin{definition}[Boss Fight]\label{def:bossfight}
Fix a hidden secret $C\in\C$ (destination/key/profile) and a declared repetition window (units, reset rule, and allowed query budget).
The \emph{Boss Fight} is the repeated-use identification game in which an analyst observes a horizon-$T$ sequence of views $Y_{1:T}$ (transcripts or projections) generated by interacting with the mechanism over that window, possibly adaptively (choosing queries, inducing retries, or selecting stopping times), and attempts to identify $C$.
\end{definition}

\paragraph{Imports.}
Threat/window declarations are imported from the series checklist note.\cite{series:threatwindows}
Mechanism-specific contracts (scheduling/stop-time padding, TV certificates, and spectral delegation bounds) are treated as black-box inputs.\cite{series:schedcompiler,series:knapsacktransport,series:prefixcapacities,series:manytestinglowerbounds,series:stoptimeaddendum,series:spectral}
The algebraic bridge between MaxL/PML and realized odds inflation is imported from the endpoint-metrics synthesis note.\cite{series:endpointbridge}

\paragraph{Exports.}
This paper exports MaxL/PML as repeated-use contracts, conservative step bounds for common transcript patterns, and necessary \emph{witness} conditions (``KICS'') showing when ``target+decoys'' collapses under partial observation and selective failure.
Papers~B--C compile these budgets into concrete anonymous-DHT knobs and proof-carrying certificates.

\paragraph{Artifact policy.}
This archive is source-only; supporting artifacts and tooling are available on request.\cite{series:artifactpolicy}


\section{Positioning within the series}\label{sec:series}
This paper supplies the \emph{identification-budget lens}.
Mechanism design for termination-time hiding, schedule equalization, and TV/KL padding is developed elsewhere in the series
(e.g., \cite{series:schedcompiler,series:knapsacktransport,series:prefixcapacities,series:manytestinglowerbounds,series:stoptimeaddendum}).
Random-walk delegation and spectral distinguishability bounds for overlays are treated in \cite{series:spectral}.
Here we focus on:
(i) MaxL/PML as an end-to-end \emph{repeated-use} contract,
(ii) support-local bounds suited to code review, and
(iii) necessary witness conditions (``KICS'') that prevent decoy collapse.

\section{Threat model: repeated-use identification}\label{sec:bossfight}
Let $C\in\C$ be a hidden secret (destination, key, etc.) and let one protocol instance emit an observation $Y\in\Y$ (a transcript or any analyst feature extraction).
In the Boss Fight, an adversary sees a sequence $(Y_1,\dots,Y_T)$ and may choose queries, retries, or stopping times adaptively.
The engineering question is: \emph{how fast can an optimal analyst identify $C$ as $T$ grows, and how do defenses add up?}

\paragraph{Threat/window declarations (imported).}
To keep the Boss-Fight sub-series non-redundant, we treat threat-model boilerplate as ``owned'' by the series checklist note on threat models and repetition windows~\cite{series:threatwindows}.
This paper assumes a repeated-use identification window and focuses on endpoint accounting; Papers~B--C instantiate the checklist fields for concrete DHT transcripts.

\section{MaxL/PML as sharp identification budgets}\label{sec:maxl}
We recall the operational definition; see \cite{issa2020operational}.

\begin{definition}[Maximal leakage \cite{issa2020operational}]
For discrete $C,Y$,
\[
\MaxL(C\!\to\!Y)\;\triangleq\;\logtwo\sum_{y\in\Y}\max_{c\in\C}\Prb[Y=y\given C=c].
\]
\end{definition}


\begin{definition}[Conditional maximal leakage]\label{def:condmaxl}
For a public context variable $H$ (history, adversary-chosen inputs, configuration), define the \emph{contextual} leakage as the worst-case over realized contexts:
\[
\MaxL(C\!\to\!Y\mid H)\;\triangleq\;\sup_{h\in\supp(H)}\;\logtwo\sum_{y\in\Y}\max_{c\in\C}\Prb[Y=y\given C=c,H=h].
\]
Likewise define the conditional pointwise leakage for an outcome $y$ under context $h$ by
\[
\PML(C\!\to\!y\mid h)\;\triangleq\;\logtwo\frac{\max_{c}\Prb[Y=y\given C=c,H=h]}{\Prb[Y=y\given H=h]}.
\]
\end{definition}
\paragraph{Operational sizing rule (imported).}
MaxL/PML admit a clean ``exact identification'' sizing rule: for a uniform $K$-way secret, keeping exact-guess probability below $\alpha$ over a window requires total MaxL at most $\log_2(\alpha K)$ bits.
We treat the algebraic bridge between MaxL/PML and realized odds inflation as owned by the endpoint-metrics synthesis note~\cite{series:endpointbridge}.

\begin{table}[t]
\centering
\caption{Total Boss Fight leakage $B$ required to keep exact-guess probability below $\varepsilon$ when $C$ is uniform on $K$ secrets, using $P_{\mathrm{guess}}(C\mid Y^{T})\le 2^{B}/K$. Values are the largest $B$ (bits) that still ensures $P_{\mathrm{guess}}\le \varepsilon$; negative values mean even $B=0$ is insufficient.}
\label{tab:identification-targets}
\begin{tabular}{r r r r}
\toprule
$K$ & $\varepsilon=0.01$ & $\varepsilon=0.001$ & $\varepsilon=1e-06$\\
\midrule
$2^{16}$ & 9.36 & 6.03 & -3.93 \\
$2^{20}$ & 13.36 & 10.03 & 0.07 \\
$2^{24}$ & 17.36 & 14.03 & 4.07 \\
\bottomrule
\end{tabular}
\end{table}


\begin{table}[t]
\centering
\caption{A minimal ``Boss Fight'' sizing calculation. Even when the secret space is only $K=2^{20}$ and the analyst's exact-guess probability target is $\delta=10^{-6}$, the total allowable MaxL budget in the window is $B\le \logtwo(\delta K)\approx 0.068$ bits. If a client performs $Q=1000$ lookups in the window, the average per-lookup budget is $b=B/Q\approx 6.8\times 10^{-5}$ bits.}
\label{tab:window-example}
\begin{tabular}{rrrr}
\toprule
$K$ & $\delta$ & $B\le \logtwo(\delta K)$ (bits) & $b=B/Q$ for $Q{=}1000$ \\
\midrule
$2^{20}$ & $10^{-6}$ & $0.068$ & $6.8\times 10^{-5}$ \\
\bottomrule
\end{tabular}
\end{table}




\paragraph{Unconditional PML.}
We write $\PML(C\!\to\!y)$ for Definition~\ref{def:condmaxl}'s pointwise leakage when the context is empty (or absorbed into $Y$). Its operational meaning---worst-case multiplicative inflation of posterior odds for \emph{any} secret event---is given in \cite{saeidian2023pml,saeidian2025rethinking}.

\begin{lemma}[Data processing for MaxL/PML \cite{issa2020operational,saeidian2023pml}]\label{lem:dp}
If $Z=f(Y)$ is any (possibly randomized) post-processing of $Y$, then
$\MaxL(C\!\to\!Z)\le \MaxL(C\!\to\!Y)$.
Moreover, for every outcome $z$, $\PML(C\!\to\!z)\le \max_{y\in\supp(Y\mid Z=z)} \PML(C\!\to\!y)$ (pointwise leakage cannot increase when outcomes are merged).
\end{lemma}

\begin{proof}[Proof sketch]
For MaxL, for any channel $P_{Z|Y}$,
\[
\sum_z \max_c P(z\mid c)
=\sum_z \max_c \sum_y P(z\mid y)P(y\mid c)
\le \sum_z \sum_y P(z\mid y)\max_c P(y\mid c)
= \sum_y \max_c P(y\mid c),
\]
using $\max_c\sum_y a_y(c)\le \sum_y \max_c a_y(c)$ and $\sum_z P(z\mid y)=1$.
Taking $\logtwo$ gives $\MaxL(C\!\to\!Z)\le \MaxL(C\!\to\!Y)$.
For PML, bound the likelihood ratio at $z$ by a convex combination of likelihood ratios over $y$ in the preimage, hence it is at most the maximum preimage ratio.
\end{proof}


\begin{lemma}[Adaptive composition accountant]\label{lem:adaptive}
Let $(Y_1,\dots,Y_T)$ be produced sequentially, where at round $t$ the mechanism may depend on the public history $H_{t-1}\triangleq (Y_1,\dots,Y_{t-1})$ (and any adversary-chosen inputs) but the hidden secret is a fixed $C$.
Then
\[
\MaxL(C\!\to\!Y_{1:T}) \;\le\; \sum_{t=1}^{T}\MaxL(C\!\to\!Y_t \mid H_{t-1}),
\]
and for every realized transcript $y_{1:T}$,
\[
\PML(C\!\to\!y_{1:T}) \;\le\; \sum_{t=1}^{T}\PML(C\!\to\!y_t \mid y_{1:t-1}).
\]
\end{lemma}
\begin{proof}[Proof sketch]
For pointwise leakage, expand the likelihood ratio:
$\frac{\max_c p(y_{1:T}\mid c)}{p(y_{1:T})}
=\frac{\max_c \prod_t p(y_t\mid c,y_{<t})}{\prod_t p(y_t\mid y_{<t})}
\le \prod_t \frac{\max_c p(y_t\mid c,y_{<t})}{p(y_t\mid y_{<t})}$,
and take $\logtwo$.
For MaxL, an additive adaptive composition bound follows from data processing and the sequential characterization of maximal leakage (with the contextual definition above); see \cite{issa2020operational,issa2022adaptive,saeidian2023pml}.
\end{proof}

\begin{corollary}[Boss Fight firewall: budget $\rightarrow$ repetition cap]\label{cor:firewall}
Suppose $C$ is uniform over $|\C|$ secrets and each protocol instance $t$ satisfies the contextual bound $\MaxL(C\!\to\!Y_t\mid H_{t-1})\le b$.
Then after $T$ instances,
\[
P_{\mathrm{guess}}(C\mid Y_{1:T})\le \frac{2^{Tb}}{|\C|}.
\]
Thus, to keep $P_{\mathrm{guess}}(C\mid Y_{1:T})\le \delta$, it suffices to enforce the conservative cap
\[
T \le \left\lfloor \frac{\logtwo(\delta|\C|)}{b}\right\rfloor.
\]
\end{corollary}

\paragraph{Remark.}
This yields a simple \emph{rate-limit / retry-cap} policy derived from a proved MaxL upper bound.
It is often pessimistic (it ignores that many rounds may reveal little), motivating the pointwise \emph{filter/odometer} alternative below.
Paper~B uses this cap as the first line of a practical dial sheet.

\paragraph{Discounts under mixing.}
The additive accountant is worst-case.
If the secret influence contracts across rounds (e.g., due to cover mixing or random-walk delegation), one can replace linear growth by a discounted sum with contraction factors; see \cite{series:spectral} for the overlay-specific derivations used in this series. Recent work also connects bounded pointwise maximal leakage to Dobrushin contraction coefficients beyond the LDP regime for discrete kernels, yielding explicit contraction factors under mild support assumptions on the relevant priors \cite{grosse2026dobrushin}.

\paragraph{Filters and odometers (operational governance).}
In repeated-use deployments, operators may want a runtime rule that \emph{stops} or \emph{throttles} once too much leakage has accumulated, rather than assuming a fixed horizon $T$.
The pointwise accountant in Lemma~\ref{lem:adaptive} makes this straightforward: maintain a running sum of realized conditional PML increments and halt when it exceeds a policy threshold (a MaxL/PML analogue of privacy filters/odometers in differential privacy) \cite{rogers2016privacyodometers}.

\paragraph{Relation to other information measures.}
Mutual information can severely understate worst-case \emph{identification} risk.
For example, let $C$ be uniform over $K$ secrets and define the ``spike'' channel $Y$ by: with probability $1-\varepsilon$, output $\bot$; with probability $\varepsilon$, output $C$.
Then $I(C;Y)=\varepsilon\logtwo K$ (tiny when $\varepsilon$ is small), but
\[
\MaxL(C\!\to\!Y)=\logtwo\big((1-\varepsilon)+\varepsilon K\big)=\logtwo\big(1+\varepsilon(K-1)\big),
\]
so even $\varepsilon\approx 1/K$ yields \emph{order-1} MaxL.
This is exactly the Boss Fight phenomenon: rare ``perfect-witness'' events accumulate under repetition.
MaxL equals Sibson mutual information of order $\infty$ and admits equivalent variational characterizations in terms of $\Dinfty$ divergences; see \cite{issa2016sibson,issa2020operational,esposito2019renyi}.
When the adversary has explicit side information $Z$, a conservative way to account for it is to use conditional variants (e.g., maximal conditional $\alpha$-leakage at $\alpha=\infty$), and to rely on robustness results showing that side information cannot increase leakage when it is conditionally independent of the release given the secret \cite{liao2019sideinfo,liao2020malprops}.
Finally, if the design goal is to protect a \emph{particular} known secret rather than an unknown function of $C$, secret-specific variants such as statistic maximal leakage can be relevant \cite{wang2024statmaxl}.

\paragraph{Remark.} Additive adaptive composition for maximal leakage is formalized in several settings; see e.g.\ \cite{issa2022adaptive} (binary inputs) and the PML composition results in \cite{saeidian2023pml}.



\paragraph{Broader context.}
Maximal leakage has been argued to better capture \emph{worst-case} adversarial success in side-channel settings than mutual information or capacity, and to yield tractable cost--leakage trade-offs as linear programs~\cite{wu2020sidechannel}.
In this series we use the direct step-wise (support-local) bounds below, because they are easy to witness and check.

\section{A support-local step bound used by the series}\label{sec:supportlocal}
Many overlays implement the pattern:
a \emph{public} cover policy (routing kernel, baseline schedule) plus \emph{secret-dependent injections} (forced destinations, aborts, terminations, or extra messages).
The series uses the following conservative bound as a checker-friendly rule.

\begin{lemma}[Forced destination + $w$ dummies (full visibility)]\label{lem:forced_dest_full}
Let $C$ be uniform on $K$ destinations.
A mechanism outputs a set $S$ of size $w{+}1$ consisting of the forced destination $C$ plus $w$ dummies sampled uniformly without replacement from the remaining $K{-}1$ destinations.
If the analyst observes $S$ fully, then
\[
\MaxL(C\!\to\!S)=\logtwo\frac{K}{w+1}.
\]
\end{lemma}

\begin{lemma}[Erasure/partial visibility rule]\label{lem:erasure}
Let $Y$ be any channel from $C$ and define the \emph{erased} output $\tilde Y$ by:
with probability $1-p$, output $\bot$ (independent of $C$), and with probability $p$, output $Y$.
Then
\[
2^{\MaxL(C\to \tilde Y)} = (1-p) + p\cdot 2^{\MaxL(C\to Y)}.
\]
Equivalently, $\MaxL(C\to \tilde Y)=\logtwo\!\big((1-p)+p\cdot 2^{\MaxL(C\to Y)}\big)$.
\end{lemma}

\begin{proof}[Proof sketch]
For $yeq\bot$, $\Prb[\tilde Y=y\mid C=c]=p\,\Prb[Y=y\mid C=c]$, and $\Prb[\tilde Y=\bot\mid C=c]=1-p$. Therefore,
\[
2^{\MaxL(C\to \tilde Y)}=\sum_{\tilde y}\max_c \Prb[\tilde Y=\tilde y\mid C=c]
=(1-p)+p\sum_{y}\max_c \Prb[Y=y\mid C=c]
=(1-p)+p\cdot 2^{\MaxL(C\to Y)}.
\]
Taking $\logtwo$ yields the stated form.
\end{proof}

Combining Lemmas~\ref{lem:forced_dest_full}--\ref{lem:erasure} gives a simple visibility-adjusted bound:
\[
\MaxL(C\!\to\!\tilde S)=\logtwo\!\Big((1-p)+p\cdot \frac{K}{w+1}\Big),
\]
which is close to $\logtwo\!\big(1+p(K-1)/w\big)$ when $w$ is large.
This is the ``support-local'' rule used in certificate payloads: the checker only needs $(K,w,p)$ (or the local surrogate $(M,w,p)$ when $K$ is replaced by a support-size summary $M$).

\section{KICS: why ``target + dummies'' collapses under partial observation}\label{sec:kics}
The rule above already implies a harsh scaling under full visibility: to keep $\MaxL$ below $b$ bits over $K$ destinations requires $w \gtrsim K/2^b$.
Partial observability and selective failure often make things worse in practice because \emph{sparse} destination witnesses accumulate under repetition.

\subsection{A closed form under partial identifier visibility}
Let $C\in\{1,\dots,N\}$ be uniform.
A mechanism outputs $S=\{C\}\cup D$ where $D$ is a uniform $k$-subset of $\{1,\dots,N\}\setminus\{C\}$.
An observer sees each element of $S$ independently with probability $q$ and sees nothing outside $S$.
Let $Y$ be the observed subset.

\begin{lemma}[Exact MaxL for partial observability]\label{lem:kics-partial-obs}
In the above model,
\[
\MaxL(C\!\to\!Y)
= \logtwo\!\left(
(1-q)^{k+1} + \sum_{m=1}^{k+1} \binom{N}{m}\frac{\binom{N-m}{k-m+1}}{\binom{N-1}{k}}\, q^{m}(1-q)^{k+1-m}
\right).
\]
Moreover, for every $m\ge 1$ the maximizer in $\max_c \Prb[Y=O\mid C=c]$ is achieved by choosing $c\in O$.
\end{lemma}
\begin{proof}[Proof sketch]
Condition on $C=c$.
For any fixed observed subset $O$ of size $m\ge 1$, the event $\{Y=O\}$ requires that all elements of $O$ are in $S$ and are observed, and all other elements of $S$ are erased.
If $c\notin O$ then $\Prb[Y=O\mid C=c]=0$.
If $c\in O$, by symmetry of the uniform dummy choice, $\Prb[Y=O\mid C=c]$ depends only on $m$ and equals
$q^{m}(1-q)^{k+1-m}\cdot \binom{N-m}{k-m+1}/\binom{N-1}{k}$.
Summing $\max_c \Prb[Y=O\mid C=c]$ over all subsets $O$ grouped by size $m$ yields the closed form inside the $\logtwo$.
\end{proof}



\begin{table}[t]
\centering
\caption{Partial-visibility collapse for ``target+$k$ dummies'' over $N=2^{16}$ destinations with per-identifier visibility $q=0.2$. We report $\MaxL(C\!\to\!Y)$ in bits using the exact closed form in Lemma~\ref{lem:kics-partial-obs}. Even thousands of dummies leave multi-bit leakage; getting below $1$ bit requires $k\approx 0.5N$, and getting below $0.25$ bits requires $k>0.8N$.}
\label{tab:kics-collapse}
\begin{tabular}{r r}
\toprule
$k$ & $\MaxL$ (bits)\\
\midrule
0 & 13.678\\
5 & 12.976\\
10 & 12.411\\
20 & 11.594\\
50 & 10.328\\
100 & 9.342\\
500 & 7.031\\
1000 & 6.033\\
5000 & 3.712\\
10000 & 2.712\\
20000 & 1.712\\
40000 & 0.712\\
60000 & 0.127\\
\bottomrule
\end{tabular}
\end{table}

\paragraph{Interpretation.}
Let $S=\{C\}\cup D$ be the full set of size $k{+}1$ (full visibility).
Then the partial-visibility observation $Y$ is obtained from $S$ by independent erasure of each element with probability $1-q$, hence by Lemma~\ref{lem:dp},
$\MaxL(C\!\to\!Y)\le \MaxL(C\!\to\!S)=\logtwo\!\big(N/(k{+}1)\big)$.
Table~\ref{tab:kics-collapse} shows that this upper bound is essentially \emph{achieved} for realistic $q$:
partial visibility helps only when $q$ is extremely small (Table~\ref{tab:visibility-sweep}).
Thus ``target+$k$ dummies'' is not a scalable knob in the Boss Fight regime: to drive per-lookup leakage below a bit, $k$ must be a large fraction of the destination universe.


\begin{table}[t]
\centering
\caption{Visibility erasure helps only when per-identifier visibility is extremely rare. We fix $N=2^{16}$, $k=1000$ dummies, and compute $\MaxL(C\!\to\!Y)$ under the partial-visibility model of Lemma~\ref{lem:kics-partial-obs} for different per-identifier visibility $q$. For $q\gtrsim 0.05$, $\MaxL$ is essentially the full-visibility value $\logtwo(N/(k{+}1))\approx 6.03$ bits.}
\label{tab:visibility-sweep}
\begin{tabular}{r r}
\toprule
$q$ & $\MaxL$ (bits)\\
\midrule
$1\times 10^{-6}$ & 0.090\\
$1\times 10^{-4}$ & 2.836\\
$1\times 10^{-3}$ & 5.385\\
0.01 & 6.033\\
0.05 & 6.033\\
0.10 & 6.033\\
0.20 & 6.033\\
0.50 & 6.033\\
0.99 & 6.033\\
\bottomrule
\end{tabular}
\end{table}

\subsection{KICS as necessary witness constraints}
We use ``KICS'' (key indistinguishability covering constraints) as shorthand for a hierarchy of constraints on destination-indexed witnesses:
marginal (1-wise), pairwise (2-wise), and $\ell$-wise witness events.
The Boss Fight takeaway is: under repetition, marginal cover is typically insufficient; safety requires higher-order KICS or mechanisms whose abort/timeout behavior is provably post-processing.

\section{What to import from earlier papers}\label{sec:import}
For building concrete defenses (deadline bucketing, stop-time padding, poset-feasible schedules, and TV/KL certificates), we refer to earlier papers in this series
\cite{series:schedcompiler,series:knapsacktransport,series:prefixcapacities,series:stoptimeaddendum}.
In Papers~B--C we show how these imported mechanisms are used under a MaxL/PML accountant and how to ship checkable certificates.

\subsection{A minimal dictionary: from series certificates to MaxL/PML}
Earlier papers in this series often certify privacy via \emph{TV/KL-style} constraints on transcript distributions (for termination-time hiding, schedule equalization, and padding kernels).
To use those mechanisms in the Boss Fight regime, we need \emph{identification} budgets (MaxL/PML).
The intended interface is to extract from the earlier certificate artifacts a \emph{public baseline} distribution $Q$ (the compiler's cover kernel) and a small set of \emph{witnessed envelopes} bounding the secret-conditioned mass relative to $Q$.

\begin{lemma}[Envelope-to-MaxL interface]\label{lem:envelope}
Let $Q$ be a public distribution on $\Y$ and suppose the implementation (or its certificate) guarantees that for all secrets $c$ and outputs $y$,
\[
\Prb[Y=y\mid C=c]\;\le\; R\cdot Q(y)
\qquad\text{for some constant }R\ge 1.
\]
Then $\MaxL(C\to Y)\le \logtwo R$.
Moreover, if there exists a constant $s\in(0,1]$ such that the marginal satisfies $\Prb[Y=y]\ge s\,Q(y)$ for all $y$ with $Q(y)>0$, then for every such $y$,
\[
\PML(C\!\to\!y)\le \logtwo\!\big(R/s\big).
\]
\end{lemma}
\begin{proof}
By definition,
$\sum_y\max_c \Prb[Y=y\mid C=c]\le \sum_y RQ(y)=R$, hence $\MaxL\le\logtwo R$.
For PML, $\max_c \Prb[y\mid c]/\Prb[y]\le \max_c \Prb[y\mid c]/(sQ(y))\le R/s$.
\end{proof}


\paragraph{Where does $R$ come from?}
In the series, $R$ is typically witnessed by \emph{bucketed} or \emph{deadline} constraints (e.g., per-bucket mixture bounds) and is checkable from dual certificates (see the stop-time and max-mixture addenda \cite{series:stoptimeaddendum} and the TV/knapsack certificate paper \cite{series:knapsacktransport}).
This paper does not re-derive those envelopes; it only states the Boss Fight accounting interface.
For readers who prefer an $\ell_1$ privacy contract, total-variation privacy admits a linear-program utility tradeoff and satisfies post-processing/linkage properties \cite{rassouli2018tv}.
Our stance is that MaxL/PML are the sharper budgets when the adversary's objective is \emph{identification under repetition}.

\paragraph{Which contract when?}
\begin{itemize}
\item \textbf{MaxL / PML (this paper):} when the attacker is trying to \emph{identify} a secret (or a function of it) and can \emph{repeat} or \emph{adapt} queries; MaxL/PML are operationally equal to worst-case multiplicative gain in guessing and admit clean composition \cite{issa2020operational,saeidian2023pml}.
\item \textbf{TV / KL (earlier papers):} when the goal is distributional \emph{indistinguishability} of transcripts (simulation-style arguments, padding-kernel design, stop-time equalization). TV contracts are convex/LP-friendly and satisfy post-processing/linkage \cite{rassouli2018tv}.
\item \textbf{Local DP:} when the desired guarantee is an explicit \emph{per-outcome likelihood-ratio bound} (useful under arbitrary side information), at the cost of often looser identification budgets; see standard LDP mechanism structure \cite{duchi2013localprivacy,kairouz2016ldp}.
\end{itemize}

\section{Takeaways}
\begin{itemize}
\item \textbf{Budgets should match the adversary's goal:} if the attacker repeats lookups to \emph{identify} a destination, MaxL/PML are the clean endpoint metrics.
\item \textbf{MaxL composes cleanly:} per-epoch identification budgets add across epochs, enabling ``privacy-to-knobs'' compilation and switchable policy menus.
\item \textbf{KICS matters:} ``target + dummies'' can collapse under projections/partial visibility; treat witness constraints (and post-processing) as first-class and certify them.
\end{itemize}

\section*{Acknowledgments}
This draft is part of an anonymity series. Artifacts and certificates are available on request.

\begin{thebibliography}{99}
\bibitem{issa2020operational}
I.~Issa, A.~B. Wagner, and S.~Kamath.
\newblock An operational approach to information leakage.
\newblock \emph{IEEE Transactions on Information Theory}, 66(3):1625--1657, 2020.
\newblock DOI:10.1109/TIT.2019.2962804.
\newblock Open preprint: arXiv:1807.07878. \url{https://arxiv.org/abs/1807.07878}.

\bibitem{issa2016sibson}
I.~Issa, S.~Kamath, and A.~B. Wagner.
\newblock Maximal leakage minimization for the Shannon cipher system.
\newblock ISIT 2016; extended manuscript PDF. \url{https://sudeepkamath.github.io/files/maximal_leakage_minimization_IKW.pdf}.

\bibitem{issa2022adaptive}
I.~Issa and A.~B. Wagner.
\newblock An adaptive composition theorem for maximal leakage for binary inputs.
\newblock In \emph{Proc. IEEE International Symposium on Information Theory (ISIT)}, 2022.
\newblock DOI:10.1109/ISIT50566.2022.9834392.

\bibitem{saeidian2023pml}
S.~Saeidian, G.~Cervia, T.~J. Oechtering, and M.~Skoglund.
\newblock Pointwise maximal leakage.
\newblock \emph{IEEE Transactions on Information Theory}, 69(12):8054--8080, 2023.
\newblock DOI:10.1109/TIT.2023.3304378.

\bibitem{saeidian2025rethinking}
S.~Saeidian, G.~Cervia, T.~J. Oechtering, and M.~Skoglund.
\newblock Rethinking disclosure prevention with pointwise maximal leakage.
\newblock \emph{Journal of Privacy and Confidentiality}, 2025.
\newblock Open PDF: \url{https://journalprivacyconfidentiality.org/index.php/jpc/article/download/893/780/2658}.
\newblock Open preprint: arXiv:2303.07782. \url{https://arxiv.org/abs/2303.07782}.

\bibitem{series:endpointbridge}
Anonymous.
\newblock Endpoint Metrics Bridge for Anonymous Overlays: Odds Inflation, Pointwise Maximal Leakage, and Maximal Leakage.
\newblock Series synthesis note, 2026. (This archive.)

\bibitem{series:threatwindows}
Anonymous.
\newblock Threat Models and Repetition Windows for Anonymous DHTs: a short declaration checklist.
\newblock Series synthesis note, 2026. (This archive.)

\bibitem{esposito2019renyi}
A.~R. Esposito, M.~Gastpar, and I.~Issa.
\newblock Generalization error bounds via R\'enyi-, $f$-divergences and maximal leakage.
\newblock arXiv:1912.01439, 2019.

\bibitem{liao2019sideinfo}
J.~Liao, L.~Sankar, O.~Kosut, and F.~P. Calmon.
\newblock Robustness of maximal $\alpha$-leakage to side information.
\newblock arXiv:1901.07105, 2019.

\bibitem{liao2020malprops}
J.~Liao, L.~Sankar, O.~Kosut, and F.~P. Calmon.
\newblock Maximal $\alpha$-leakage and its properties.
\newblock 2020. PDF: \url{https://par.nsf.gov/servlets/purl/10232244}.

\bibitem{wang2024statmaxl}
S.~Wang, Z.~Lin, and G.~Fanti.
\newblock Statistic maximal leakage.
\newblock arXiv:2411.18531, 2024.

\bibitem{wu2020sidechannel}
B.~Wu, A.~B. Wagner, and G.~E. Suh.
\newblock A case for maximal leakage as a side channel leakage metric.
\newblock arXiv:2004.08035, 2020.

\bibitem{duchi2013localprivacy}
J.~C. Duchi, M.~I. Jordan, and M.~J. Wainwright.
\newblock Local privacy, data processing inequalities, and statistical minimax rates.
\newblock arXiv:1302.3203, 2013.
\newblock (See also the NIPS 2014 version.)

\bibitem{kairouz2016ldp}
P.~Kairouz, S.~Oh, and P.~Viswanath.
\newblock Extremal mechanisms for local differential privacy.
\newblock \emph{Journal of Machine Learning Research}, 17(1):1--51, 2016.
\newblock \url{https://jmlr.org/papers/v17/15-135.html}.

\bibitem{grosse2025empiricalpml}
L.~Grosse, S.~Saeidian, M.~Skoglund, and T.~J.~Oechtering.
\newblock Privacy mechanism design based on empirical distributions.
\newblock arXiv:2509.22428, 2025.

\bibitem{grosse2026dobrushin}
L.~Grosse, S.~Saeidian, T.~J. Oechtering, and M.~Skoglund.
\newblock Dobrushin coefficients of private mechanisms beyond local differential privacy.
\newblock arXiv:2601.09498, 2026. \url{https://arxiv.org/abs/2601.09498}.


\bibitem{rassouli2018tv}
B.~Rassouli and D.~G\"und\"uz.
\newblock Optimal utility-privacy trade-off with the total variation distance as the privacy measure.
\newblock \emph{IEEE Transactions on Information Forensics and Security}, 15, 2020.
\newblock DOI:10.1109/TIFS.2019.2903658.
\newblock Open preprint: arXiv:1801.02505.

\bibitem{rogers2016privacyodometers}
R.~M. Rogers, A.~Roth, J.~Ullman, and S.~Vadhan.
\newblock Privacy odometers and filters: pay-as-you-go composition.
\newblock In \emph{Advances in Neural Information Processing Systems (NeurIPS)}, 2016.
\newblock arXiv:1605.08294. \url{https://arxiv.org/abs/1605.08294}.

\bibitem{series:schedcompiler}
Anonymous.
\newblock \emph{Scheduling and Compiler Techniques for Metadata-Hiding Overlay Lookups}.
\newblock Series paper (published), 2026.
\newblock Wiki: \texttt{[[2026.01.22 - Mathematics: Scheduling and Compiler Techniques for Metadata-Hiding Overlay Lookups]]}.


\bibitem{series:knapsacktransport}
Anonymous.
\newblock \emph{Optimal Poset-Feasible Padding: Knapsack, Transport, and Dual Certificates for TV Privacy}.
\newblock Series paper (published), 2026.
\newblock Wiki: \texttt{[[2026.01.29 - Mathematics: Optimal Poset-Feasible Padding: Knapsack, Transport, and Dual Certificates for TV Privacy]]}.


\bibitem{series:prefixcapacities}
Anonymous.
\newblock \emph{Prefix Capacities and Separation Cuts for Shared-Kernel Termination-Time Hiding}.
\newblock Series paper (published), 2026.
\newblock Wiki: \texttt{[[2026.01.25 - Mathematics: Prefix Capacities and Separation Cuts for Shared-Kernel Termination-Time Hiding]]}.


\bibitem{series:manytestinglowerbounds}
Anonymous.
\newblock \emph{Many-Testing Lower Bounds for Low-Latency Privacy on Hierarchical Overlays}.
\newblock Series paper (published), 2026. (Optional deep-dive in this archive: \texttt{published/optional\_many\_testing\_lower\_bounds/paper.tex}.)


\bibitem{series:stoptimeaddendum}
Anonymous.
\newblock \emph{Stop-Time Padding Addendum (unique-only): Max-Mixture Separations and Tail-Sign Deadline Normal Forms}.
\newblock Series note (published), 2026.
\newblock Wiki: \texttt{[[2026.01.26 - Mathematics: Stop-Time Padding Addendum, Max-Mixture Separations and Tail-Sign Deadline Normal Forms]]}.


\bibitem{series:spectral}
Anonymous.
\newblock \emph{Spectral Anonymity and Optimal Distinguishing Bounds for Random-Walk Delegation in (Possibly Directed) Overlays}.
\newblock Series paper (published), 2026.
\newblock Wiki: \texttt{[[2026.01.23 - Mathematics: Spectral Anonymity and Optimal Distinguishing Bounds for Random-Walk Delegation in (Possibly Directed) Overlays]]}.


\bibitem{series:artifactpolicy}
Anonymous.
\newblock Artifact policy (source-only archive; artifacts available on request).
\newblock Series note, 2026.

\bibitem{series:notation}
Anonymous.
\newblock Notation and Minimal Model for the One-True-Archive.
\newblock Series synthesis note (Synthesis~8), 2026. (This archive.)


\end{thebibliography}

\end{document}